\documentclass[final]{beamer}

% Compile on Overleaf with pdfLaTeX or XeLaTeX.
% Poster size: A0 portrait.
\usepackage[size=a0,orientation=portrait,scale=1.18]{beamerposter}
\usepackage[T1]{fontenc}
\usepackage[utf8]{inputenc}
\usepackage[polish]{babel}
\usepackage{graphicx}
\usepackage{ragged2e}
\usepackage{tikz}
\usepackage[most]{tcolorbox}
\usepackage{qrcode}
\usepackage{fontawesome5}
\usepackage[sfdefault]{quicksand}

\usetikzlibrary{positioning,shapes.geometric}

\definecolor{DeepTeal}{HTML}{0D4A48}
\definecolor{MainTeal}{HTML}{5FB7B9}
\definecolor{DarkTeal}{HTML}{0B7975}
\definecolor{SoftAqua}{HTML}{A9DADA}
\definecolor{PaleAqua}{HTML}{EEF8F8}
\definecolor{MutedTeal}{HTML}{4C7372}
\definecolor{AccentViolet}{HTML}{9B5CFF}
\definecolor{WarmYellow}{HTML}{FFF7DF}

\setbeamercolor{normal text}{fg=DeepTeal,bg=white}
\setbeamercolor{block title}{fg=DeepTeal,bg=PaleAqua}
\setbeamercolor{block body}{fg=DeepTeal,bg=white}
\setbeamertemplate{navigation symbols}{}
\setlength{\parindent}{0pt}
\setlength{\parskip}{0.5em}

\newcommand{\posterurl}{https://dermato-xai.vercel.app/}
\newcommand{\projecttitle}{DermatoskopiaAI}
\newcommand{\projectsubtitle}{Wyjaśnialne demo adnotacji struktur dermoskopowych}

\tcbset{
  medicalcard/.style={
    enhanced,
    colback=white,
    colframe=SoftAqua,
    coltitle=DeepTeal,
    arc=14mm,
    boxrule=1.5pt,
    left=18mm,
    right=18mm,
    top=14mm,
    bottom=14mm,
    drop fuzzy shadow={black!10!white},
    fonttitle=\bfseries\Large,
  },
  tealcard/.style={
    enhanced,
    colback=MainTeal,
    colframe=MainTeal,
    coltext=white,
    arc=18mm,
    boxrule=0pt,
    left=22mm,
    right=22mm,
    top=18mm,
    bottom=18mm,
    drop fuzzy shadow={MainTeal!35!white},
  },
  pill/.style={
    enhanced,
    colback=white,
    colframe=SoftAqua,
    arc=10mm,
    boxrule=1pt,
    left=10mm,
    right=10mm,
    top=5mm,
    bottom=5mm,
  }
}

\newcommand{\sectioncard}[2]{
  \begin{tcolorbox}[medicalcard,title={#1}]
    \justifying
    #2
  \end{tcolorbox}
}

\newcommand{\metricbox}[3]{
  \begin{tcolorbox}[
    enhanced,
    colback=PaleAqua,
    colframe=PaleAqua,
    arc=10mm,
    boxrule=0pt,
    width=\linewidth,
    left=10mm,
    right=10mm,
    top=8mm,
    bottom=8mm
  ]
    {\Huge\bfseries\color{DarkTeal}#1}\\[0.2em]
    {\Large\bfseries #2}\\
    {\normalsize\color{MutedTeal}#3}
  \end{tcolorbox}
}

\begin{document}
\begin{frame}[t]
\begin{tikzpicture}[remember picture,overlay]
  \fill[SoftAqua!55]
    ([xshift=-8cm,yshift=5cm]current page.north west)
    ellipse [x radius=18cm, y radius=10cm, rotate=-18];
  \fill[SoftAqua!45]
    ([xshift=8cm,yshift=-5cm]current page.south east)
    ellipse [x radius=20cm, y radius=12cm, rotate=-18];
  \fill[AccentViolet!13]
    ([xshift=-2cm,yshift=-18cm]current page.north east)
    circle [radius=5.5cm];
\end{tikzpicture}

\vspace*{1.0cm}
\begin{columns}[T,totalwidth=\textwidth]
  \begin{column}{0.12\textwidth}
    \begin{tcolorbox}[
      enhanced,
      colback=MainTeal,
      colframe=MainTeal,
      arc=14mm,
      boxrule=0pt,
      width=\linewidth,
      height=6.5cm,
      valign=center,
      halign=center
    ]
      {\Huge\bfseries\color{white}DAI}
    \end{tcolorbox}
  \end{column}
  \begin{column}{0.68\textwidth}
    {\fontsize{72}{78}\selectfont\bfseries\color{DeepTeal}\projecttitle}\\[0.15em]
    {\fontsize{34}{40}\selectfont\bfseries\color{DarkTeal}\projectsubtitle}\\[0.25em]
    {\Large\color{MutedTeal}AI Forum 2026 \quad | \quad Publiczna wizytówka i demo projektu}
  \end{column}
  \begin{column}{0.16\textwidth}
    \begin{tcolorbox}[medicalcard,width=\linewidth,halign=center]
      {\Large\bfseries Demo}\\[0.6em]
      \qrcode[height=5.2cm]{\posterurl}\\[0.6em]
      {\small\color{MutedTeal}\texttt{\posterurl}}
    \end{tcolorbox}
  \end{column}
\end{columns}

\vspace{1cm}

\begin{tcolorbox}[tealcard,width=\textwidth]
  {\fontsize{36}{42}\selectfont\bfseries
  Jak człowiek i model AI różnią się w oznaczaniu struktur dermoskopowych?}\\[0.35em]
  {\Large
  Plakat pokazuje działający przepływ: użytkownik zaznacza struktury na obrazie,
  model RF-DETR + SAHI wykonuje detekcję, a aplikacja porównuje obie adnotacje
  przez nakładkę, statystyki i edukacyjny feedback. Narzędzie ma charakter
  badawczo-edukacyjny i nie zastępuje konsultacji lekarskiej.}
\end{tcolorbox}

\vspace{1cm}

\begin{columns}[T,totalwidth=\textwidth]
  \begin{column}{0.32\textwidth}
    \sectioncard{Problem}{
      Dermoskopia wymaga rozpoznawania subtelnych struktur wizualnych.
      W projekcie koncentrujemy się na klasach takich jak yellow globules,
      blue-gray globules, rosettes, milia-like cyst i MAY globules.

      \vspace{0.4em}
      \textbf{Wyzwanie:} samo pokazanie predykcji modelu nie wyjaśnia, gdzie
      człowiek zgadza się z modelem, a gdzie pojawia się różnica.
    }

    \sectioncard{Cel demo}{
      \begin{itemize}
        \item umożliwić szybkie zaznaczenie struktur bez logowania;
        \item porównać adnotacje użytkownika z detekcjami modelu;
        \item pokazać wynik w formie zrozumiałej na plakacie i telefonie;
        \item jasno oddzielić demonstrację od zastosowania diagnostycznego.
      \end{itemize}
    }

    \sectioncard{Dane i współpraca}{
      Demo powstało dzięki współpracy klinicznej i wsparciu przy oznaczaniu
      obrazów dermoskopowych.

      \vspace{0.4em}
      \textbf{Partner kliniczny:} Old Town Clinic\\
      \textbf{Wsparcie merytoryczne:} Dr hab. n. med. Jacek Calik,
      Lek. Karolina Gowans\\
      \textbf{Opieka naukowa i techniczna:} Grzegorz Chodak,
      Przemysław Dolata, Piotr Giedziun
    }
  \end{column}

  \begin{column}{0.34\textwidth}
    \sectioncard{Przepływ użytkownika}{
      \begin{enumerate}
        \item Użytkownik wybiera obraz przykładowy albo wgrywa własny.
        \item Wybiera aktywne klasy struktur.
        \item Zaznacza bounding boxy przez kliknięcie dwóch rogów.
        \item Aplikacja porównuje adnotacje z predykcjami modelu.
        \item Wynik pokazuje nakładkę, statystyki i feedback tekstowy.
      \end{enumerate}
    }

    \begin{tcolorbox}[medicalcard,title={Miejsce na zrzut ekranu demo}]
      \begin{tikzpicture}
        \draw[rounded corners=18pt, fill=PaleAqua, draw=SoftAqua, line width=2pt]
          (0,0) rectangle (\linewidth,18cm);
        \node[align=center, text=MutedTeal, font=\Large\bfseries]
          at (0.5\linewidth,9cm)
          {Wstaw tutaj screenshot aplikacji\\
          \texttt{\textbackslash includegraphics[width=\textbackslash linewidth]\{demo-screenshot.png\}}};
      \end{tikzpicture}
    \end{tcolorbox}

    \sectioncard{Model i porównanie}{
      Model detekcji działa jako zewnętrzny endpoint Modal.
      Konfiguracja obejmuje wieloklasową detekcję z RF-DETR + SAHI,
      osobnymi progami pewności oraz łączeniem predykcji.

      \vspace{0.4em}
      Porównanie wykorzystuje IoU dla bounding boxów i rozdziela wynik na:
      dopasowania, adnotacje użytkownika bez zgodnej predykcji oraz predykcje
      modelu bez zgodnej adnotacji użytkownika.
    }
  \end{column}

  \begin{column}{0.30\textwidth}
    \sectioncard{Co pokazujemy?}{
      \metricbox{5}{klas struktur}{konfigurowalny zestaw widoczny dla użytkownika}
      \vspace{0.5em}
      \metricbox{0}{logowania}{publiczny dostęp z QR kodu}
      \vspace{0.5em}
      \metricbox{1}{interaktywny przepływ}{upload, adnotacja, porównanie, feedback}
    }

    \sectioncard{Najważniejsze elementy XAI}{
      \begin{itemize}
        \item \textbf{Nakładka wizualna:} linie użytkownika i modelu są rozdzielone.
        \item \textbf{Statystyki:} liczba dopasowań i rozbieżności jest jawna.
        \item \textbf{Feedback:} opisuje różnice językiem edukacyjnym.
        \item \textbf{Kontrola klas:} użytkownik wybiera, które struktury są analizowane.
      \end{itemize}
    }

    \sectioncard{Ograniczenia}{
      \begin{itemize}
        \item Demo nie jest narzędziem diagnostycznym.
        \item Wynik zależy od jakości obrazu i wybranych klas.
        \item Porównanie bounding boxów upraszcza złożoną ocenę kliniczną.
        \item Feedback tekstowy ma charakter edukacyjny, nie medyczny.
      \end{itemize}
    }

    \sectioncard{Zespół}{
      Marcel Masiukiewicz \quad
      Szymon Szafraniec \quad
      Bartłomiej Ruszaj \quad
      Marcin Tkocz

      \vspace{0.5em}
      {\color{MutedTeal}Projekt przygotowany jako publiczna prezentacja i demo
      funkcjonalności na AI Forum 2026.}
    }
  \end{column}
\end{columns}

\vfill
\begin{tcolorbox}[
  enhanced,
  colback=PaleAqua,
  colframe=PaleAqua,
  arc=12mm,
  boxrule=0pt,
  width=\textwidth,
  left=16mm,
  right=16mm,
  top=8mm,
  bottom=8mm
]
  \begin{columns}[T,totalwidth=\textwidth]
    \begin{column}{0.74\textwidth}
      {\Large\bfseries\color{DeepTeal}Przetestuj demo na telefonie}\\
      {\color{MutedTeal}
      Zeskanuj QR kod, wybierz przykładowe zdjęcie i porównaj swoje adnotacje
      z detekcjami modelu.}
    \end{column}
    \begin{column}{0.22\textwidth}
      \RaggedLeft
      {\Large\bfseries\color{DarkTeal}\faIcon{globe}\ \texttt{\posterurl}}
    \end{column}
  \end{columns}
\end{tcolorbox}

\end{frame}
\end{document}
